\chapter*{General Conclusion}
\addcontentsline{toc}{chapter}{General Conclusion}
\label{ch:general_conclusion}

The primary objective of this Final Year Project was to design, develop, and deploy a comprehensive digital platform dedicated to the learning and perfection of Quranic recitation and Tajweed rules. Through the development of \textbf{Tahkik}, we have demonstrated how modern technology --- specifically Artificial Intelligence, cross-platform mobile development, and cloud-based backend services --- can be leveraged to address the traditional challenges of accessibility and personalized feedback in Islamic education.

In Chapter~\ref{ch:state_of_art}, we presented the domain of Quranic recitation and Tajweed, analyzed existing digital solutions including Quran.com, Tarteel AI, and online tutoring platforms, and identified a critical gap: the absence of a platform combining instant AI feedback with structured human teacher review.

In Chapter~\ref{ch:conceptual}, we modeled the system using UML to ensure a clear separation of concerns among the different actors (Learners, Teachers, Administrators), defined 10 functional and 5 non-functional requirements, and designed the system architecture with detailed use case, sequence, and class diagrams.

In Chapter~\ref{ch:implementation}, we successfully implemented the platform using Flutter for the mobile app, React for the web dashboard, Supabase for the backend, and Python with the Whisper model for AI-powered recitation analysis.

\begin{emeraldbox}[title=Key Contributions]
\begin{enumerate}
  \item A cross-platform mobile application providing an intuitive Quran reading and recitation recording experience.
  \item Integration of the Whisper AI model for real-time pronunciation analysis and Tajweed error detection.
  \item A dedicated web dashboard empowering teachers to review, correct, and manage student recitations.
  \item A comprehensive administration panel with analytics, user management, and AI health monitoring.
  \item A human-in-the-loop architecture combining AI efficiency with teacher expertise.
\end{enumerate}
\end{emeraldbox}

Despite these achievements, certain limitations remain:
\begin{itemize}
  \item The AI model's accuracy varies across different recitation styles (Qira'at) and regional accents.
  \item The current implementation requires an internet connection for AI processing.
  \item The gamification features and social interaction elements are limited in this version.
\end{itemize}

\begin{goldbox}[title=Future Perspectives]
\begin{itemize}
  \item Fine-tuning the Whisper model specifically on diverse Qira'at (Warsh, Hafs, etc.) for improved accuracy.
  \item Adding offline AI inference capabilities using on-device models.
  \item Implementing gamification features (badges, leaderboards, learning paths) to increase engagement.
  \item Integrating live tutoring sessions with built-in video/audio calls between teachers and students.
  \item Developing community features for learners to share achievements and support each other.
\end{itemize}
\end{goldbox}

\begin{emeraldbox}[title=Closing Statement]
\begin{center}
\large\itshape
``Read! In the name of your Lord who created.''\\[6pt]
\normalsize\normalfont\color{black} --- Surah Al-Alaq, Verse 1
\end{center}
\vspace{6pt}
\normalsize
The Tahkik platform represents a meaningful step forward in digital Islamic education. By harmonizing the efficiency of AI with the irreplaceable expertise of human teachers, Tahkik aspires to make high-quality Tajweed instruction accessible to a global audience, preserving the oral tradition of the Quran for generations to come.
\end{emeraldbox}
